\documentclass[a4paper,11pt]{article}

\usepackage[utf8]{inputenc}
\usepackage[T1]{fontenc}
\usepackage[margin=2.5cm]{geometry}
\usepackage{amsmath,amssymb}
\usepackage{listings}
\usepackage{color}
\usepackage{hyperref}
\usepackage{parskip}
\usepackage{enumitem}
\usepackage{tabularx}
\usepackage{xcolor}
\usepackage{titlesec}

\definecolor{codegreen}{rgb}{0,0.6,0}
\definecolor{codegray}{rgb}{0.5,0.5,0.5}
\definecolor{codepurple}{rgb}{0.58,0,0.82}
\definecolor{backcolour}{rgb}{0.95,0.95,0.92}

\lstdefinestyle{mystyle}{
    backgroundcolor=\color{backcolour},
    commentstyle=\color{codegreen},
    keywordstyle=\color{magenta},
    numberstyle=\tiny\color{codegray},
    stringstyle=\color{codepurple},
    basicstyle=\ttfamily\footnotesize,
    breakatwhitespace=false,
    breaklines=true,
    captionpos=b,
    keepspaces=true,
    numbers=none,
    numbersep=5pt,
    showspaces=false,
    showstringspaces=false,
    showtabs=false,
    tabsize=2,
    language=C,
    frame=single,
}
\lstset{style=mystyle}

\hypersetup{
    colorlinks=true,
    linkcolor=blue,
    urlcolor=blue,
    citecolor=blue
}

\title{Maxxwell Server Interface (MXSI) \\ \large Server Programming Guide}
\author{The Maxxwell Kernel Team}
\date{Version 1.0 -- July 2026}

\begin{document}

\maketitle
\thispagestyle{empty}
\newpage

\tableofcontents
\newpage

\section{Introduction}

The Maxxwell Server Interface (MXSI) is a server management framework
built into the Maxxwell hosted microkernel. It provides the
infrastructure for userspace servers to register, communicate, and
be managed by the kernel throughout their lifecycle.

MXSI treats every userspace service as a {\em server} --- an
independent process with a well-defined identity, a set of
capabilities, and a message-based communication channel managed by
the kernel. This design enables strong isolation, structured
service discovery, and centralized lifecycle control.

This document serves as a guide for developers who wish to write new
servers for the Maxxwell operating system. It covers the
architecture, the C API exposed by MXSI, and presents a complete
worked example.

\subsection{Core Concepts}

\begin{description}[leftmargin=2cm,style=nextline]
  \item[Server] A named, registered userspace process with a unique
    ID, a state machine, and one or more Mach ports for communication.
  \item[Registry] A kernel-maintained directory of all live servers,
    indexed by ID, name, and capability.
  \item[Capability] A bitmask that describes what a server can do
    (e.g., filesystem, networking, display). Clients discover servers
    by querying capabilities.
  \item[Message] A structured kernel-mediated datagram that carries
    a type, flags, a sender port, a destination port, and a
    variable-length payload.
  \item[Manager] The kernel subsystem responsible for starting,
    stopping, and restarting servers.
  \item[Bootstrap] The boot-time sequence that brings up the first
    server and, if that fails, falls back to the MsWaiter kernel
    monitor.
\end{description}

\section{Architecture Overview}

\subsection{Layers}

MXSI is composed of six subsystems, each exposed through a header
in \texttt{Kernel/mxwl/server/}:

\begin{center}
\begin{tabularx}{\textwidth}{|l|X|}
\hline
\textbf{Module} & \textbf{Purpose} \\
\hline
\texttt{server.h/c}   & Core server object --- data structures,
                        state machine, initialisation and teardown. \\
\texttt{registry.h/c} & Server directory --- insertion, removal,
                        lookup by ID, name, or capability. \\
\texttt{manager.h/c}  & Lifecycle control --- start, stop,
                        restart, poll, MsWaiter command dispatch. \\
\texttt{message.h/c}  & Message definition and manipulation ---
                        create, validate, duplicate, destroy. \\
\texttt{ipc.h/c}      & IPC transport layer --- send, reply,
                        broadcast, notify, RPC, endpoint management. \\
\texttt{loader.h/c}   & Server executable loader --- load, unload,
                        start, stop, IPC setup, readiness wait. \\
\texttt{bootstrap.h/c}& Boot sequence --- init, configure, start
                        default servers, MsWaiter fallback. \\
\hline
\end{tabularx}
\end{center}

\subsection{Kernel Mediation}

All inter-server communication passes through the kernel. When server
A sends a message to server B, the kernel:

\begin{enumerate}
  \item Validates the message header and payload.
  \item Verifies that server B is in the \texttt{READY} or
    \texttt{WAITING} state.
  \item Resolves B's endpoint port from the registry.
  \item Delivers the message to B's receive queue.
  \item Optionally invokes a global monitor callback.
\end{enumerate}

This design gives the kernel full visibility into all IPC traffic,
enabling debugging, tracing, and access control.

\begin{center}
\texttt{\small
  Server A -\hbox{-}$>${\bf kernel} -\hbox{-}$>$ Server B
}
\end{center}

\section{Server Lifecycle}

\subsection{State Machine}

Every server transitions through the following states:

\begin{center}
\begin{tabularx}{\textwidth}{|c|X|}
\hline
\textbf{State} & \textbf{Meaning} \\
\hline
\texttt{STOPPED}  & Not loaded or initialised. \\
\texttt{LOADING}  & Server executable is being loaded into memory. \\
\texttt{STARTING} & Server process has been spawned. \\
\texttt{READY}    & Server is running and accepting requests. \\
\texttt{WAITING}  & Server is idle, waiting for external events. \\
\texttt{FAILED}   & Server failed during start-up. \\
\texttt{PANIC}    & Server encountered an unrecoverable error. \\
\texttt{STOPPING} & Server is being shut down. \\
\hline
\end{tabularx}
\end{center}

\subsection{Initialisation}

Every server begins as a \texttt{mx\_server\_t} object allocated by
the kernel or the bootstrap code. Initialisation sets the name,
assigns a unique numeric ID, and resets all fields:

\begin{lstlisting}
#include <mxsi.h>

mx_server_t *my_server = libkern_calloc(1, sizeof(mx_server_t));
kern_return_t kr = mx_server_init(my_server, "my-server");
if (kr != KERN_SUCCESS) {
    /* handle error */
}
\end{lstlisting}

After \texttt{mx\_server\_init}, the server is in the
\texttt{STOPPED} state. The server's capabilities, flags, and path
should be set before it is registered:

\begin{lstlisting}
my_server->capabilities = MX_CAP_PROCESS | MX_CAP_IPC;
my_server->flags        = MX_SERVER_SYSTEM | MX_SERVER_AUTOSTART;
libkern_strncpy(my_server->path, "/mach/servers/my-server",
                MX_SERVER_PATH_MAX - 1);
\end{lstlisting}

\subsection{Registration}

Before a server can be managed or addressed, it must be registered
with the kernel's registry:

\begin{lstlisting}
kern_return_t kr = mx_server_register(my_server);
if (kr == KERN_NAME_EXISTS) {
    /* A server with this name already exists */
}
\end{lstlisting}

\subsection{Teardown}

When a server is removed from the system, it should be unregistered
and destroyed:

\begin{lstlisting}
mx_server_unregister(my_server);
mx_server_destroy(my_server);
libkern_free(my_server);
\end{lstlisting}

\section{The Registry}

The registry is a kernel-maintained directory of all servers. It
supports four kinds of lookup:

\subsection{Lookup by ID}

\begin{lstlisting}
mx_server_t *srv = mx_server_lookup_id(42);
\end{lstlisting}

\subsection{Lookup by Name}

\begin{lstlisting}
mx_server_t *srv = mx_server_lookup_name("4bsd");
\end{lstlisting}

\subsection{Lookup by Capability}

Returns the first server that has the exact capability:

\begin{lstlisting}
mx_server_t *srv = mx_server_lookup_capability(MX_CAP_FILESYSTEM);
\end{lstlisting}

\subsection{Lookup by Capability Mask}

Returns the first server that has {\em any} of the given
capabilities:

\begin{lstlisting}
mx_server_t *srv = mx_server_lookup_capability_any(
    MX_CAP_DISPLAY | MX_CAP_INPUT);
\end{lstlisting}

\subsection{Fuzzy Search}

The \texttt{mx\_server\_find} function accepts a query string that
may match a server by name or ID (parsed as an integer):

\begin{lstlisting}
mx_server_t *srv;
kern_return_t kr = mx_server_find("4bsd", &srv);
\end{lstlisting}

\subsection{Enumeration}

The registry supports iteration over all servers:

\begin{lstlisting}
mx_server_t *s = mx_server_first();
while (s) {
    printf("server: %s\n", s->name);
    s = mx_server_next(s);
}
\end{lstlisting}

To get a snapshot of server info into a caller-supplied buffer:

\begin{lstlisting}
mx_server_info_t buf[64];
int n = mx_server_list(buf, 64);
\end{lstlisting}

\subsection{Constants}

\begin{lstlisting}
#define MX_SERVER_MAX     64  /* maximum registered servers */
#define MX_SERVER_NAME_MAX    64
#define MX_SERVER_VERSION_MAX 32
#define MX_SERVER_PATH_MAX    256
\end{lstlisting}

\section{Messages}

All inter-server communication is carried by \texttt{mx\_message\_t}
objects. A message has a fixed header and a variable-length payload
stored in a flexible array member.

\subsection{Message Structure}

\begin{lstlisting}
typedef struct mx_message {
    mach_port_t   sender;        /* port of the sending server  */
    mach_port_t   destination;   /* port of the target server   */
    uint32_t      msg_id;        /* application-defined ID      */
    mx_msg_type_t type;          /* SYNC, ASYNC, NOTIFY, etc.   */
    uint32_t      payload_size;  /* size of payload in bytes    */
    uint32_t      flags;         /* MX_MSG_* flags              */
    uint8_t       payload[];     /* variable-length data        */
} mx_message_t;
\end{lstlisting}

\subsection{Message Types}

\begin{center}
\begin{tabularx}{\textwidth}{|c|X|}
\hline
\textbf{Type} & \textbf{Semantics} \\
\hline
\texttt{MX\_MSG\_SYNC}    & Sender blocks waiting for a reply.\\
\texttt{MX\_MSG\_ASYNC}   & Fire-and-forget: no reply is expected.\\
\texttt{MX\_MSG\_NOTIFY}  & One-way event notification.\\
\texttt{MX\_MSG\_REPLY}   & Response to a previous message.\\
\texttt{MX\_MSG\_REQUEST} & A typed request (application-level).\\
\texttt{MX\_MSG\_RESPONSE}& A typed response.\\
\hline
\end{tabularx}
\end{center}

\subsection{Message Flags}

\begin{lstlisting}
MX_MSG_NONE         = 0
MX_MSG_URGENT       = (1 << 0)   /* high-priority delivery */
MX_MSG_EXPECT_REPLY = (1 << 1)   /* sender expects a reply */
MX_MSG_REPLY        = (1 << 2)   /* this is itself a reply */
MX_MSG_BROADCAST    = (1 << 3)   /* intended for broadcast */
MX_MSG_NOTIFY       = (1 << 4)   /* notification message   */
\end{lstlisting}

\subsection{Creating Messages}

\begin{lstlisting}
mx_message_t *msg;
kern_return_t kr = mx_msg_create(
    &msg,               /* output pointer          */
    sender_port,        /* source port             */
    dest_port,          /* destination port        */
    0x1000,             /* message ID              */
    MX_MSG_SYNC,        /* type                    */
    MX_MSG_NONE,        /* flags                   */
    data,               /* payload pointer         */
    data_size           /* payload size            */
);
\end{lstlisting}

\subsection{Destroying Messages}

\begin{lstlisting}
mx_msg_destroy(msg);
\end{lstlisting}

\subsection{Validation and Duplication}

\begin{lstlisting}
kern_return_t kr = mx_msg_validate(msg);
mx_message_t *copy = mx_msg_dup(msg);
\end{lstlisting}

\subsection{Utility}

\begin{lstlisting}
uint32_t total = mx_msg_total_size(msg);
\end{lstlisting}

Returns the sum of the header size and the payload size.

\section{IPC Transport Layer}

The IPC module provides the functions that actually deliver messages
between servers. It operates on Mach ports that are allocated during
server endpoint creation.

\subsection{Endpoint Management}

Each server has two ports:
\begin{itemize}
  \item \textbf{endpoint} --- the service port where clients send
    requests.
  \item \textbf{mgmt\_port} --- the management port used by the
    kernel for lifecycle control (stop, restart, etc.).
\end{itemize}

\begin{lstlisting}
kern_return_t kr = mx_ipc_server_endpoint_create(my_server);
if (kr == KERN_SUCCESS) {
    /* my_server->endpoint and my_server->mgmt_port are now valid */
}
\end{lstlisting}

Destroy endpoints when the server is unloaded:

\begin{lstlisting}
mx_ipc_server_endpoint_destroy(my_server);
\end{lstlisting}

\subsection{Sending Messages}

\subsubsection{Unicast Send}

Send a message to a specific server:

\begin{lstlisting}
kern_return_t kr = mx_server_send(target_server, msg);
\end{lstlisting}

The kernel checks that the target is in \texttt{READY} or
\texttt{WAITING} state and that its endpoint port is valid.

\subsubsection{Reply}

Reply to a previously received request:

\begin{lstlisting}
kern_return_t kr = mx_server_reply(request_msg, reply_msg);
\end{lstlisting}

The reply's destination is automatically set to the original
sender's port.

\subsubsection{Broadcast}

Send a message to all servers whose capabilities match a given mask:

\begin{lstlisting}
kern_return_t kr = mx_server_broadcast(
    MX_CAP_FILESYSTEM | MX_CAP_STORAGE, msg);
\end{lstlisting}

\subsubsection{Notification}

Send a one-way event notification to a server:

\begin{lstlisting}
kern_return_t kr = mx_server_notify(target, event_id,
                                    data, data_size);
\end{lstlisting}

\subsubsection{Synchronous RPC}

Send a request and block until a reply arrives:

\begin{lstlisting}
mx_message_t *reply = NULL;
kern_return_t kr = mx_server_rpc(target, request,
                                 &reply, timeout_ms);
if (kr == KERN_SUCCESS) {
    /* process reply */
    mx_msg_destroy(reply);
}
\end{lstlisting}

\subsection{Polling}

The kernel calls \texttt{mx\_ipc\_poll()} to drain pending messages
from all server endpoints. Servers that integrate with the kernel's
event loop do not need to call this directly; the manager invokes it
from \texttt{mx\_manager\_poll()}.

\subsection{Monitoring}

A global callback can be installed to observe every message that
passes through the IPC layer:

\begin{lstlisting}
void my_monitor(mx_message_t *msg, void *ctx) {
    klog_info("msg %d sent\n", msg->msg_id);
}

mx_ipc_set_monitor(my_monitor, NULL);
\end{lstlisting}

This is used by MsWaiter's tracing feature.

\section{Server Manager}

The manager is the kernel subsystem that controls the server
lifecycle. It ties together the registry, loader, and IPC layer.

\subsection{Lifecycle Functions}

\begin{lstlisting}
/* Start a server by its registry ID or name */
kern_return_t mx_server_start(server_id);
kern_return_t mx_server_start_by_name("my-server");

/* Stop a running server */
kern_return_t mx_server_stop(server_id);
kern_return_t mx_server_stop_by_name("my-server");

/* Restart (stop then start) */
kern_return_t mx_server_restart(server_id);
kern_return_t mx_server_restart_by_name("my-server");
\end{lstlisting}

\subsection{Direct Server Operations}

When a \texttt{mx\_server\_t*} pointer is already available:

\begin{lstlisting}
kern_return_t kr = mx_manager_start_server(server);
kern_return_t kr = mx_manager_stop_server(server);
kern_return_t kr = mx_manager_restart_server(server);
\end{lstlisting}

\subsection{Initialisation and Shutdown}

\begin{lstlisting}
kern_return_t kr = mx_manager_init();
/* ... */
mx_manager_shutdown();
\end{lstlisting}

\subsection{Command Dispatch}

The manager provides a command dispatch function used by the
MsWaiter kernel monitor. It accepts argc/argv arrays and routes
them to internal command handlers:

\begin{lstlisting}
char out[4096];
kern_return_t kr = mx_manager_handle_command(
    2, (char *[]){"start", "4bsd", NULL}, out, sizeof(out));
printf("%s\n", out);
\end{lstlisting}

Supported commands:

\begin{center}
\begin{tabularx}{\textwidth}{|l|X|}
\hline
\textbf{Command} & \textbf{Description} \\
\hline
\texttt{servers}   & List all registered servers. \\
\texttt{start}     & Start a server by name. \\
\texttt{stop}      & Stop a server by name. \\
\texttt{restart}   & Restart a server by name. \\
\texttt{info}      & Show detailed information about a server. \\
\texttt{trace}     & Enable/disable IPC tracing for a server. \\
\texttt{panic}     & Manually set a server to PANIC state. \\
\hline
\end{tabularx}
\end{center}

\section{Loader}

The loader subsystem is responsible for loading a server executable,
creating the host process (via fork/exec on POSIX), setting up IPC
endpoints, and waiting for the server to signal readiness.

\subsection{Load and Unload}

\begin{lstlisting}
kern_return_t kr = mx_loader_load(server, "/path/to/server");
kern_return_t kr = mx_loader_unload(server);
\end{lstlisting}

\subsection{Start and Stop}

\begin{lstlisting}
kern_return_t kr = mx_loader_start(server);
kern_return_t kr = mx_loader_stop(server);
\end{lstlisting}

\subsection{IPC Setup and Readiness}

\begin{lstlisting}
kern_return_t kr = mx_loader_setup_ipc(server);
kern_return_t kr = mx_loader_wait_ready(server, timeout_ms);
\end{lstlisting}

The readiness wait blocks until the server sends a \texttt{READY}
notification on its management port, or until the timeout expires.

\section{Capabilities}

Capabilities are a bitmask that servers advertise at registration
time. They allow clients to discover the right server for a given
task without knowing its name or ID.

\subsection{Capability Constants}

\begin{lstlisting}
MX_CAP_NONE       = 0
MX_CAP_FILESYSTEM = (1 << 0)   /* filesystem services      */
MX_CAP_DISPLAY    = (1 << 1)   /* display / framebuffer    */
MX_CAP_NETWORK    = (1 << 2)   /* networking stack         */
MX_CAP_AUDIO      = (1 << 3)   /* audio services           */
MX_CAP_PROCESS    = (1 << 4)   /* process management       */
MX_CAP_BSD        = (1 << 5)   /* BSD compatibility layer  */
MX_CAP_DEVICE     = (1 << 6)   /* device I/O access        */
MX_CAP_STORAGE    = (1 << 7)   /* block storage            */
MX_CAP_IPC        = (1 << 8)   /* IPC forwarding services  */
MX_CAP_MONITOR    = (1 << 9)   /* system monitoring        */
MX_CAP_INPUT      = (1 << 10)  /* input handling           */
MX_CAP_BOOTSTRAP  = (1U << 31) /* bootstrap server         */
\end{lstlisting}

\subsection{Capability String}

Convert a capability bitmask to a human-readable comma-separated
string:

\begin{lstlisting}
char buf[MX_SERVER_CAPS_STR_MAX];
const char *str = mx_server_caps_string(server->capabilities,
                                         buf, sizeof(buf));
\end{lstlisting}

\section{Bootstrap Sequence}

The bootstrap subsystem orchestrates the initial server start-up
when the kernel boots.

\subsection{Boot Flow}

\begin{enumerate}
  \item The kernel initialises: IPC subsystem $\rightarrow$ tasks
    $\rightarrow$ loader $\rightarrow$ MsWaiter monitor.
  \item \texttt{mx\_bootstrap\_init()} initialises the bootstrap
    configuration with a 30-second timeout and autostart enabled.
  \item If a boot argument \texttt{-srv=/path/to/server} was
    provided, that path is saved.
  \item \texttt{mx\_bootstrap\_start()} initialises the registry,
    the IPC layer, the loader, and the manager in sequence.
  \item \texttt{mx\_bootstrap\_sequence()} runs:
    \begin{enumerate}
      \item If a configured server path exists, it attempts to load
        and start that server as the bootstrap server with
        \texttt{MX\_CAP\_BOOTSTRAP | MX\_CAP\_PROCESS | MX\_CAP\_IPC}
        capabilities.
      \item If the configured server fails or no path was given,
        \texttt{mx\_bootstrap\_launch\_mswaiter()} is called,
        which enters the kernel-resident MsWaiter monitor.
    \end{enumerate}
\end{enumerate}

\subsection{API Reference}

\begin{lstlisting}
kern_return_t mx_bootstrap_init(void);
kern_return_t mx_bootstrap_configure(const char *boot_args);

kern_return_t mx_bootstrap_start(void);

kern_return_t mx_bootstrap_start_server(
    const char *name,
    const char *path,
    uint32_t capabilities);

kern_return_t mx_bootstrap_launch_mswaiter(void);
kern_return_t mx_bootstrap_sequence(void);

const char  *mx_bootstrap_get_server_path(void);
\end{lstlisting}

\subsection{Adding a Server to the Bootstrap}

Servers can be added to the boot sequence by calling
\texttt{mx\_bootstrap\_start\_server} from a custom bootstrap
implementation:

\begin{lstlisting}
mx_bootstrap_start_server(
    "my-server",
    "/mach/servers/my-server",
    MX_CAP_FILESYSTEM | MX_CAP_STORAGE);
\end{lstlisting}

\section{Complete Example: A Minimal Server}

This section walks through a complete server that registers with the
kernel, handles a simple ping request, and shuts down gracefully. It
is intended as a template for new server authors.

\subsection{Source Code}

\begin{lstlisting}
/* my_server.c -- minimal MXSI server */

#include <mxsi.h>
#include <stdio.h>
#include <string.h>
#include <unistd.h>

#define MY_SERVER_NAME   "my-server"
#define MY_SERVER_VERSION "1.0"
#define MSG_PING         0x1000
#define MSG_SHUTDOWN     0x1001

/* ------------------------------------------------------------------ */
/*  Message handler                                                    */
/* ------------------------------------------------------------------ */

static kern_return_t handle_ping(mx_message_t *msg,
                                  mx_message_t **reply)
{
    return mx_msg_create(reply, MACH_PORT_NULL, msg->sender,
                         0x2000, MX_MSG_REPLY, MX_MSG_NONE,
                         "pong", 5);
}

static kern_return_t handle_shutdown(mx_message_t *msg,
                                      mx_message_t **reply)
{
    (void)msg;
    (void)reply;
    exit(0);
    return KERN_SUCCESS;
}

/* ------------------------------------------------------------------ */
/*  Dispatch loop                                                      */
/* ------------------------------------------------------------------ */

static void dispatch_loop(mx_server_t *server)
{
    while (1) {
        /* Receive a message on our endpoint */
        mx_message_t msg;
        kern_return_t kr = mach_msg_receive(
            server->endpoint, (mach_msg_t *)&msg,
            sizeof(msg) + MX_MSG_MAX_PAYLOAD, 1000);

        if (kr != KERN_SUCCESS)
            continue;

        mx_message_t *reply = NULL;

        switch (msg.msg_id) {
        case MSG_PING:
            handle_ping(&msg, &reply);
            break;
        case MSG_SHUTDOWN:
            handle_shutdown(&msg, &reply);
            break;
        default:
            /* unknown message -- ignore */
            break;
        }

        if (reply) {
            mx_server_reply(&msg, reply);
            mx_msg_destroy(reply);
        }
    }
}

/* ------------------------------------------------------------------ */
/*  Main entry point                                                   */
/* ------------------------------------------------------------------ */

int main(int argc, char *argv[])
{
    (void)argc;
    (void)argv;

    /* ---- 1. Create and initialise the server object ---- */
    mx_server_t *server = calloc(1, sizeof(mx_server_t));
    kern_return_t kr = mx_server_init(server, MY_SERVER_NAME);
    if (kr != KERN_SUCCESS) {
        fprintf(stderr, "mx_server_init failed: %d\n", kr);
        return 1;
    }

    libkern_strncpy(server->version, MY_SERVER_VERSION,
                    MX_SERVER_VERSION_MAX - 1);
    server->capabilities = MX_CAP_NONE;
    server->flags = MX_SERVER_SYSTEM;

    /* ---- 2. Register with the kernel ---- */
    kr = mx_server_register(server);
    if (kr != KERN_SUCCESS) {
        fprintf(stderr, "mx_server_register failed: %d\n", kr);
        mx_server_destroy(server);
        free(server);
        return 1;
    }

    /* ---- 3. Set up IPC endpoints ---- */
    kr = mx_ipc_server_endpoint_create(server);
    if (kr != KERN_SUCCESS) {
        fprintf(stderr, "endpoint creation failed: %d\n", kr);
        mx_server_unregister(server);
        mx_server_destroy(server);
        free(server);
        return 1;
    }

    /* ---- 4. Signal READY to the kernel ---- */
    mx_server_set_state(server, MX_SERVER_READY);

    mx_message_t *ready_msg = NULL;
    mx_msg_create(&ready_msg, MACH_PORT_NULL,
                  server->mgmt_port, 0x80000001,
                  MX_MSG_NOTIFY, MX_MSG_NONE, NULL, 0);
    if (ready_msg) {
        mach_msg_send(server->mgmt_port,
                      (mach_msg_t *)ready_msg,
                      mx_msg_total_size(ready_msg), 1000);
        mx_msg_destroy(ready_msg);
    }

    fprintf(stderr, "%s v%s READY on endpoint %d\n",
            MY_SERVER_NAME, MY_SERVER_VERSION, server->endpoint);

    /* ---- 5. Enter the dispatch loop ---- */
    dispatch_loop(server);

    /* ---- 6. Cleanup (unreachable in this example) ---- */
    mx_ipc_server_endpoint_destroy(server);
    mx_server_unregister(server);
    mx_server_destroy(server);
    free(server);
    return 0;
}
\end{lstlisting}

\subsection{Explanation}

\begin{enumerate}
  \item The server allocates and initialises a
    \texttt{mx\_server\_t} object. The name must be unique across
    all running servers.
  \item It registers with the kernel so that other servers and the
    manager can discover it.
  \item IPC endpoints are created --- these are Mach ports backed by
    Unix domain socket pairs. The endpoint port is the public
    service port; the mgmt port is reserved for kernel management
    messages (stop, restart, etc.).
  \item The server sets its state to \texttt{READY} and sends a
    notification to the kernel on its management port. The kernel
    waits for this notification in \texttt{mx\_loader\_wait\_ready}
    before considering the server fully started.
  \item The dispatch loop calls \texttt{mach\_msg\_receive} to block
    on the endpoint port, then dispatches messages to handler
    functions. Handlers may produce reply messages that are sent
    back to the original sender.
\end{enumerate}

\subsection{Building}

To build a server, include \texttt{mxsi.h} and link with the MXSI
library:

\begin{lstlisting}[language=bash]
cc -I/path/to/mxwl -c my_server.c -o my_server.o
cc -o my_server my_server.o -lmxsi -lpthread
\end{lstlisting}

\subsection{Adding to the Bootstrap}

Servers can be added to the boot sequence from a custom bootstrap
module, or started manually from the MsWaiter prompt:

\begin{lstlisting}[language=bash]
mxwl> start /mach/servers/my-server
\end{lstlisting}

\section{API Reference}

\subsection{mxwl/server.h --- Core Server Object}

\begin{lstlisting}
kern_return_t mx_server_init(mx_server_t *server, const char *name);
void          mx_server_destroy(mx_server_t *server);

kern_return_t mx_server_set_state(mx_server_t *server,
                                   mx_server_state_t state);
mx_server_state_t mx_server_get_state(const mx_server_t *server);

const char *mx_server_state_string(mx_server_state_t state);
const char *mx_server_caps_string(uint32_t caps, char *buf, size_t n);
\end{lstlisting}

\subsection{mxwl/registry.h --- Server Directory}

\begin{lstlisting}
kern_return_t  mx_registry_init(void);

kern_return_t  mx_server_register(mx_server_t *server);
kern_return_t  mx_server_unregister(mx_server_t *server);

mx_server_t   *mx_server_lookup_id(mx_server_id_t server_id);
mx_server_t   *mx_server_lookup_name(const char *name);
mx_server_t   *mx_server_lookup_capability(mx_capability_t cap);
mx_server_t   *mx_server_lookup_capability_any(uint32_t caps);

kern_return_t  mx_server_find(const char *query, mx_server_t **out);

int            mx_server_list(mx_server_info_t *buf, int max);
int            mx_server_count(void);
mx_server_t   *mx_server_first(void);
mx_server_t   *mx_server_next(mx_server_t *current);
\end{lstlisting}

\subsection{mxwl/manager.h --- Lifecycle Control}

\begin{lstlisting}
kern_return_t mx_manager_init(void);
void          mx_manager_shutdown(void);

kern_return_t mx_server_start(mx_server_id_t server_id);
kern_return_t mx_server_start_by_name(const char *name);
kern_return_t mx_server_stop(mx_server_id_t server_id);
kern_return_t mx_server_stop_by_name(const char *name);
kern_return_t mx_server_restart(mx_server_id_t server_id);
kern_return_t mx_server_restart_by_name(const char *name);

kern_return_t mx_manager_start_server(mx_server_t *server);
kern_return_t mx_manager_stop_server(mx_server_t *server);
kern_return_t mx_manager_restart_server(mx_server_t *server);

void          mx_manager_poll(void);
kern_return_t mx_manager_handle_command(int argc, char **argv,
                                         char *out, size_t out_size);
\end{lstlisting}

\subsection{mxwl/message.h --- Message Passing Primitives}

\begin{lstlisting}
typedef struct mx_message {
    mach_port_t   sender;
    mach_port_t   destination;
    uint32_t      msg_id;
    mx_msg_type_t type;
    uint32_t      payload_size;
    uint32_t      flags;
    uint8_t       payload[];
} mx_message_t;

kern_return_t mx_msg_create(mx_message_t **out, mach_port_t sender,
                             mach_port_t destination, uint32_t msg_id,
                             mx_msg_type_t type, uint32_t flags,
                             const void *payload, uint32_t payload_size);
void          mx_msg_destroy(mx_message_t *msg);
kern_return_t mx_msg_validate(const mx_message_t *msg);
mx_message_t *mx_msg_dup(const mx_message_t *msg);
uint32_t      mx_msg_total_size(const mx_message_t *msg);
\end{lstlisting}

\subsection{mxwl/ipc.h --- IPC Transport}

\begin{lstlisting}
kern_return_t mx_ipc_init(void);
void          mx_ipc_shutdown(void);

kern_return_t mx_ipc_server_endpoint_create(mx_server_t *server);
void          mx_ipc_server_endpoint_destroy(mx_server_t *server);

kern_return_t mx_server_send(mx_server_t *dest, mx_message_t *msg);
kern_return_t mx_server_reply(mx_message_t *request,
                               mx_message_t *reply);
kern_return_t mx_server_broadcast(uint32_t caps_mask,
                                   mx_message_t *msg);
kern_return_t mx_server_notify(mx_server_t *server, uint32_t event,
                                const void *data, uint32_t data_size);
kern_return_t mx_server_rpc(mx_server_t *server, mx_message_t *request,
                             mx_message_t **reply, int timeout_ms);

void mx_ipc_set_monitor(mx_ipc_monitor_t monitor, void *context);
void mx_ipc_poll(void);
\end{lstlisting}

\subsection{mxwl/loader.h --- Server Loader}

\begin{lstlisting}
kern_return_t mx_loader_init(void);
kern_return_t mx_loader_load(mx_server_t *server,
                              const char *exec_path);
kern_return_t mx_loader_unload(mx_server_t *server);
kern_return_t mx_loader_start(mx_server_t *server);
kern_return_t mx_loader_stop(mx_server_t *server);
kern_return_t mx_loader_setup_ipc(mx_server_t *server);
kern_return_t mx_loader_wait_ready(mx_server_t *server,
                                    int timeout_ms);
\end{lstlisting}

\subsection{mxwl/bootstrap.h --- Boot Sequence}

\begin{lstlisting}
kern_return_t mx_bootstrap_init(void);
kern_return_t mx_bootstrap_configure(const char *boot_args);
kern_return_t mx_bootstrap_start(void);
kern_return_t mx_bootstrap_start_server(const char *name,
                                         const char *path,
                                         uint32_t capabilities);
kern_return_t mx_bootstrap_launch_mswaiter(void);
kern_return_t mx_bootstrap_sequence(void);
const char   *mx_bootstrap_get_server_path(void);
\end{lstlisting}

\subsection{Data Structures}

\subsubsection{\texttt{mx\_server\_t}}

\begin{lstlisting}
typedef struct mx_server {
    mx_server_id_t    server_id;       /* unique numeric ID        */
    char              name[64];        /* human-readable name      */
    char              version[32];     /* server version string    */
    char              path[256];       /* executable path          */
    mx_server_state_t state;           /* current state            */
    mxwl_task_t      *task;            /* associated kernel task   */
    mxwl_thread_t    *thread;          /* associated kernel thread */
    mach_port_t       endpoint;        /* public service port      */
    mach_port_t       mgmt_port;       /* kernel management port   */
    uint32_t          capabilities;    /* capability bitmask       */
    uint32_t          flags;           /* MX_SERVER_* flags        */
    int               argc;            /* argument count           */
    char             *argv[64];        /* argument vector          */
    uint32_t          permissions;     /* access permissions       */
    int               restart_count;   /* number of restarts       */
    int               restart_delay_ms;/* ms between restarts      */
    mx_server_t      *next;            /* linked-list next pointer */
} mx_server_t;
\end{lstlisting}

\subsubsection{\texttt{mx\_server\_info\_t}}

\begin{lstlisting}
typedef struct mx_server_info {
    mx_server_id_t    server_id;
    char              name[64];
    char              version[32];
    mx_server_state_t state;
    uint32_t          capabilities;
    uint32_t          flags;
    mach_port_t       endpoint;
    int               pid;
} mx_server_info_t;
\end{lstlisting}

\section{Common Patterns}

\subsection{Server Discovery}

If a server needs to find another server at runtime, it can use the
registry lookup functions. The capability-based lookup is preferred
because it decouples the client from the server's name:

\begin{lstlisting}
mx_server_t *fs = mx_server_lookup_capability(MX_CAP_FILESYSTEM);
if (fs) {
    mx_message_t *req, *rep;
    mx_msg_create(&req, my_port, fs->endpoint, ...);
    mx_server_rpc(fs, req, &rep, 5000);
}
\end{lstlisting}

\subsection{Handling Multiple Message Types}

A common pattern is a dispatch table indexed by \texttt{msg\_id}:

\begin{lstlisting}
typedef kern_return_t (*handler_t)(mx_message_t *req,
                                    mx_message_t **rep);

typedef struct {
    uint32_t  msg_id;
    handler_t handler;
} dispatch_entry_t;

static dispatch_entry_t _handlers[] = {
    { MSG_PING,     handle_ping     },
    { MSG_LOOKUP,   handle_lookup   },
    { MSG_SHUTDOWN, handle_shutdown },
    { 0, NULL }
};
\end{lstlisting}

\subsection{Robustness}

Servers should handle the following scenarios gracefully:

\begin{itemize}
  \item \textbf{Invalid messages:} Validate \texttt{msg\_id},
    \texttt{payload\_size}, and \texttt{type} before processing.
    Return an error response or ignore the message.
  \item \textbf{Timeouts:} When calling \texttt{mx\_server\_rpc},
    always provide a reasonable timeout to avoid blocking
    indefinitely.
  \item \textbf{Shutdown:} Monitor the management port for a
    shutdown notification from the kernel. On receipt, flush
    pending work and exit cleanly.
  \item \textbf{Resource leaks:} Every \texttt{mx\_msg\_create}
    must be paired with a \texttt{mx\_msg\_destroy}. Every
    registered server should be unregistered before destruction.
\end{itemize}

\section{Appendix: Return Codes}

MXSI functions return \texttt{kern\_return\_t} values defined in
\texttt{mach/mach\_types.h}:

\begin{center}
\begin{tabularx}{\textwidth}{|c|l|X|}
\hline
\textbf{Code} & \textbf{Constant} & \textbf{Meaning} \\
\hline
0  & \texttt{KERN\_SUCCESS}           & Operation completed. \\
1  & \texttt{KERN\_INVALID\_ARGUMENT} & NULL pointer or bad parameter. \\
2  & \texttt{KERN\_FAILURE}           & General failure. \\
4  & \texttt{KERN\_INVALID\_OBJECT}   & Server or port not found. \\
5  & \texttt{KERN\_NAME\_EXISTS}      & A server with this name exists. \\
6  & \texttt{KERN\_ABORTED}           & Operation was aborted. \\
7  & \texttt{KERN\_NOT\_RECEIVER}     & Target is not ready to receive. \\
\hline
\end{tabularx}
\end{center}

\end{document}
